\section{Methods}
\label{sec:method}

\subsection{Doc-to-LoRA hypernetwork}
\label{sec:doc2lora}

\doctolora is a hypernetwork that converts a document $x$ to an LoRA adapter in one forward pass.
Specifically, a document $x$ is first sent to the base LLM, which produces a sequence of hidden states for all tokens in the document.  
These hidden states are then sent to a perceiver transformer~\citep{jaegle2021perceiver}, which compresses into a fixed-size tensor $\mathcal{U} \in \mathbb{R}^{n_L \times n_R \times n_d}$, where $n_L$ is the number of layers in the base LLM, $n_R$ is the number of rank slots, and $n_d$ is the dimension of the hidden states.
This fixed-size tensor $\mathcal{U}$ is then sent to an MLP followed by row-wise $L_2$ normalization, generating another fixed-size tensor $\mathcal{V} \in \mathbb{R}^{n_L \times n_R \times n_d}$.
Finally, a linear head $f_{\mathrm{dec},\ell}$ maps $\mathcal{V}_{\ell}$ to the LoRA adapter matrices, independently for each layer and with no weight sharing.

The resulting LoRA adapter consists of a pair of matrices $\matrix{A}_{\ell} \in \mathbb{R}^{m_\ell \times n_R}$ and $\matrix{B}_{\ell} \in \mathbb{R}^{n_R \times k_\ell}$ for each layer $\ell$, where $n_R$ is the rank of the LoRA adapter. 
These matrices $\matrix{A}_{\ell}$ and $\matrix{B}_{\ell}$ are then added to the linear layer $\matrix{W}_{\ell}$ of the MLP in the base LLM during inference.
Loading the resulting updates $\{\Delta\matrix{W}_{\ell}\}$ into the frozen base model yields a document-conditioned model that answers as if $x$ were in its context. We use the public SakanaAI checkpoint on Qwen3-4B~\citep{qwen3-2025} in the main text and report other base models in Appendix~\ref{app:benchmarks}.
The faithfulness of the decoded output depends on the hypernetwork, as evaluated by \citet{sakana2024doc2lora}.
We take this as given and focus on the embedding space itself. 
Paraphrasing the prompt changes the wording but rarely the content (Appendix~\ref{app:prompt-sensitivity}).

\subsection{\doctolora embedding}
\label{sec:doc2lora-embedding}

We leverage the post-MLP tensor $\mathcal{V}$ as the document embedding, as the following linear head $f_{\mathrm{dec},\ell}$ maps it directly to the LoRA adapter matrices, preserving linear equivariance:
\begin{equation}
f_{\mathrm{dec},\ell}\Big(\sum_i w_i\, \matrix{V}_{i,\ell}\Big) = \sum_i w_i\, f_{\mathrm{dec},\ell}(\matrix{V}_{i,\ell}).
\label{eq:linearity}
\end{equation}
Thus, averaging or interpolating embeddings is equivalent to the same operation on the adapters, a property that is lost if using the nonlinear pre-MLP matrices. Each $\mathcal{V}$ is a third-order tensor, with $\ell$-th slice $\matrix{V}_{\ell} \in \mathbb{R}^{n_R \times n_d}$ containing $L_2$-normalized, $n_d{=}512$-dimensional unit vectors per $n_R{=}8$ rank (\figref{method}).

When restricted to search, we can downsize the hidden tensor $\mathcal{V}$ while preserving much similarity information by taking the mean over the rank dimension $n_R$ (Spearman $\rho$ of $.97$--$.98$ between the full-tensor and pooled cosine similarities for Qwen3-4B, on $500$ documents from each evaluation corpus (Appendix~\ref{app:data})).
In our paper, we use the mean-pooled tensor as embedding for similarity search discussed in \secref{similarity}, while use the full tensor for generation tasks. Unless otherwise noted, we use the full tensor.

\subsection{Decoding documents and interpolants}
\label{sec:decode}

We decode an embedding by expanding $\mathcal{V}$ into a LoRA adapter and querying the frozen base model with no document in the prompt.
We illustrate with two short recipes, an Italian cacio e pepe and a Japanese kake udon, encoded with the Mistral-7B hypernetwork and decoded in full-rank mode (Appendix~\ref{app:recipe-fusion}).
With the Italian adapter loaded, the query ``What cheese is used?'' returns ``The cheese used in this recipe is Pecorino Romano.''
Next, we interpolate the two embeddings as $\mathcal{V}(\alpha)=(1-\alpha)\mathcal{V}_A+\alpha\mathcal{V}_B$, where $\mathcal{V}_A$ and $\mathcal{V}_B$ are the embeddings of the two recipes.
The endpoints $\alpha=0$ and $\alpha=1$ decode to cacio e pepe and kake udon, each with minor paraphrases.
At $\alpha=0.5$ the decode is ``Kake Udon with Creamy Parmesan Sauce.''
The dish coats udon noodles in a melted-cheese sauce and seasons them with mirin and soy.
Parmesan replaces Pecorino Romano and heavy cream appears in neither source.
We decode greedily, and the examples we report are single runs. 

\subsection{Adapting the geometry with a bijective transform}
\label{sec:kron-method}

The public hypernetwork is trained for generation.
We refine its tensor for search without retraining the base LLM and hypernetwork by adding a linear layer $g_\theta$ on top of the embedding, with a bijectivity constraint.
More specifically, for each layer $\ell$, we transform $\matrix{V}_{\ell}\in\mathbb{R}^{n_R\times n_d}$ by
\begin{align}
   \label{eq:transform}
   \matrix{V}'_{\ell} = g_\theta(\matrix{V}_{\ell}) = a_\ell \matrix{V}_{\ell} \matrix{C},
\end{align}
where $a_\ell>0$ controls the scaling of layer $\ell$, and $\matrix{C}\in\mathbb{R}^{n_d\times n_d}$ is an invertible linear matrix.
To ensure invertibility, we factor $\matrix{C}$ as $\matrix{R} \text{diag}(\sigma)$, where $\matrix{R}\in\mathbb{R}^{n_d\times n_d}$ is a rotation matrix and $\text{diag}(\sigma)\in\mathbb{R}^{n_d\times n_d}$ is a diagonal matrix with positive entries $\sigma_i>0$ (Appendix~\ref{app:bijective}).
This parametrization uses only $n_L + n_d^2$ parameters ($262$k for Qwen) and brings the embedding on par with dedicated text encoders on the similarity benchmarks we tested.
We train $g_\theta$ with InfoNCE~\citep{oord2018infonce} on citation pairs (Appendix~\ref{app:bijective}) using the OpenAlex citation pairs (negatives in-batch, as in \texttt{SPECTER2}), with the base model frozen.
